% Teil II: der durchgerechnete Weg.

\subsection*{Gegeben und gesucht}
(a) Verschiebecode $f(x)=(x+3)\bmod26$: MATHE verschlüsseln, Bijektivität
zeigen, $f^{-1}$ angeben; (b) $g(x)=2x\bmod26$ vs. $h(x)=3x\bmod26$:
welche ist bijektiv, bei der anderen zwei Kollisionen und einen ungetroffenen
Wert finden.

\subsection*{a) Verschlüsseln von MATHE}
M$=12$, A$=0$, T$=19$, H$=7$, E$=4$; Sender addiert $3$,
\begin{equation}
  \begin{array}{c|ccccc}
   &\mathrm M&\!\mathrm A&\!\mathrm T&\!\mathrm H&\!\mathrm E\\[-1pt]
   x & 12&0&19&7&4\\
   x+3\bmod26 & 15&\;3&22&10&7\\
   \text{Buchstabe} & \mathrm P&\mathrm D&\mathrm W&\mathrm K&\mathrm H
  \end{array}
\end{equation}
Also \textbf{MATHE} $\to$ \textbf{PDWKH}. Es ist eine Verschiebung, damit
bijektiv (Permutation von $\mathbb Z_{26}$); das modulare Inverse
$(+3)^{-1}=-3$:
\begin{equation}
  f^{-1}(y)=(y-3)\bmod26 .
\end{equation}

\subsection*{b) $g(x)=2x\bmod26$ und $h(x)=3x\bmod26$}
$g$: $2x$ ist stets gerade, also sind alle ungeraden Reste
($1,3,5,\dots,25$) nie erreichbar -- \emph{nicht surjektiv}, damit nicht
bijektiv. Außerdem ist $g(0)=0$ und $g(13)=0$, also $0$ und $13$ zwei
verschiedene Ursprünge mit demselben Code (nicht injektiv).

$h$: $\mathrm{ggT}(3,26)=1$, daher liefert $3x\bmod26$ alle $26$ Reste --
\emph{bijektiv}.

\begin{ergebnis}
\begin{equation*}
  \text{a)}\ \text{MATHE}\to\text{PDWKH},\quad
  f^{-1}(y)=(y-3)\bmod26,\qquad
  \text{b)}\ h\ \text{bijektiv},\ g\ \text{nicht}.
\end{equation*}
\end{ergebnis}

\begin{kontrolle}
Entschlüsseln von PDWKH mit $(y-3)\bmod26$: $15\to12=\mathrm M$,
$3\to0=\mathrm A$, $22\to19=\mathrm T$, $10\to7=\mathrm H$,
$7\to4=\mathrm E$ -- MATHE zurück. Für $g$: $1$ wird
nie getroffen; $0$ und $13$ beide $=0$. Passt.
\end{kontrolle}

\begin{merksatz}
Modulare Multiplikationen sind genau dann bijektiv, wenn der Faktor teilerfremd
zur Modulzahl ist; Verschiebungen sind immer Permutationen mit der
Umkehrung $-a$.
\end{merksatz}
